
\section{Results}\label{sec:results}

This section reports the study results by research question.  RQ1 uses the completed PyPI scan; later result tables are kept as placeholders until the remaining measurements are inserted.  The reporting structure matches the four goals in Section~\ref{sec:study} and the methodology in Section~\ref{sec:method}.

\noindent
{\bf Experiment Environment.}
Experiments are conducted in isolated Linux environments.  Package collection downloads public PyPI artifacts; later analysis phases run locally with explicit timeouts and version-specific CPython interpreters.  Runtime fuzzing is organized by CPython version, and abnormal outcomes are reported as robustness findings requiring triage rather than automatically as vulnerabilities.

\noindent
{\bf Datasets.}
The package dataset is drawn from PyPI.  For each artifact, we record distribution type, file inventory, source availability, bytecode version, hashes, and dynamic loading indicators.  RQ3 uses CPython unittest-derived bytecode seeds prepared per interpreter version rather than package bytecode, because the robustness question concerns adversarial bytecode processing rather than the benign artifacts observed in RQ1.

\noindent
{\bf Baselines and Tools.}
For RQ2 and RQ4, we use selected bytecode-analysis and decompilation tools, including standard-library loading/disassembly and available decompilers.  Tool failures are recorded with explicit reasons such as missing interpreter, timeout, unsupported bytecode, nonzero exit, or no usable source output.

\noindent
\textbf{Performance Metrics.}
We report package-level prevalence, bytecode transparency, practical analyzability levels, per-tool success and failure reasons, fuzzing outcome distributions, unique abnormal outcomes when available, and tool-bounded source-reproduction classifications.

\noindent
\textbf{Evaluation Sufficiency.}
The evaluation is structured to support the paper's motivation end to end.  RQ1 establishes whether bytecode artifacts appear and whether they are transparent to source-level inspection.  RQ2 tests whether selected tools can practically analyze the observed bytecode in a version-aware setup.  RQ3 measures the runtime boundary exposed by malformed or adversarial bytecode.  RQ4 determines whether selected source-recovery workflows can reproduce bytecode-level findings.  Together, these questions support the central claim only if prevalence, analyzability limits, and robustness findings are interpreted with their stated scope.


\subsection{RQ1: Bytecode Prevalence}\label{sec:effmeasure}

RQ1 measures how often Python bytecode artifacts appear in real distributions.  We count packages containing \texttt{.pyc} files, \texttt{\_\_pycache\_\_} directories, source-less bytecode modules, bundled bytecode artifacts, dynamically loaded code objects, and PyInstaller-style payloads.  Results should be reported both per package and per artifact, because a single package can contain many bytecode objects.

\begin{table}[t]
  \centering
  \caption{Bytecode prevalence across datasets. Replace placeholders with measured values.}
  \label{tab:prevalence}
  \begin{tabular}{lrrrr}
    \toprule
    Dataset & Packages & With \texttt{.pyc} & Source-less & Dynamic load \\
    \midrule
    Popular PyPI & TBD & TBD & TBD & TBD \\
    Random PyPI & TBD & TBD & TBD & TBD \\
    Suspicious/Malicious & TBD & TBD & TBD & TBD \\
    Bundles & TBD & TBD & TBD & TBD \\
    \bottomrule
  \end{tabular}
\end{table}

The main reporting claim for RQ1 should identify whether bytecode artifacts are rare leftovers or recurring distribution artifacts.  We also report common locations, including cache directories, package subtrees, loader resources, archive members, and bundled module stores.


\subsection{RQ2: Practical Analyzability}\label{sec:RQ2}

RQ2 asks whether observed bytecode artifacts can be analyzed by selected tools in a version-aware environment.  For each \texttt{.pyc} file found in RQ1, we run standard-library loading and disassembly, then selected decompilers when available.  Each tool result records both the status and the concrete failure reason, so failures caused by missing interpreters, unsupported bytecode versions, timeouts, and decompiler errors are not conflated.

\begin{table}[t]
  \centering
  \caption{Practical analyzability levels for observed bytecode artifacts. Replace placeholders with measured values.}
  \label{tab:tool-analyzability}
  \begin{tabular}{lrrrrr}
    \toprule
    Version & \texttt{.pyc} files & Loadable & Disassemblable & Partial decompile & Full decompile \\
    \midrule
    CPython 3.8 & TBD & TBD & TBD & TBD & TBD \\
    CPython 3.9 & TBD & TBD & TBD & TBD & TBD \\
    CPython 3.10 & TBD & TBD & TBD & TBD & TBD \\
    CPython 3.12 & TBD & TBD & TBD & TBD & TBD \\
    \bottomrule
  \end{tabular}
\end{table}

The key interpretation is practical rather than theoretical.  If bytecode becomes analyzable after the matching CPython interpreter is installed, the earlier failure is an environment limitation.  If selected tools still fail in a matched environment, the result indicates a current tooling limitation.  RQ2 therefore reports what the selected analysis stack can inspect today, not a proof that the bytecode is or is not recoverable by all possible tools.

\subsection{RQ3: Tooling Gap}\label{sec:RQ3}

RQ3 evaluates how well current tools recover or analyze Python bytecode artifacts.  For each decompiler, we measure decompilation success, syntactic validity of recovered source, recompilation success, bytecode equivalence after recompilation, and behavior preservation when safe execution is possible.  We also test whether source-level scanners observe behavior that exists only in bytecode artifacts.

\begin{table}[t]
  \centering
  \caption{Recovery workflow outcomes. Replace placeholders with measured values.}
  \label{tab:recovery}
  \begin{tabular}{lrrrr}
    \toprule
    Tool & Decompiled & Recompiled & Equivalent & Main failure \\
    \midrule
    pycdc & TBD & TBD & TBD & TBD \\
    Other & TBD & TBD & TBD & TBD \\
    \bottomrule
  \end{tabular}
\end{table}

Failure causes include unsupported Python versions, unsupported opcodes, unusual control flow, exception-table mismatches, malformed code-object metadata, and dynamic loading patterns that hide the bytecode from ordinary file-based analysis.  The main interpretation is conservative: a decompiler failure is a tooling limitation, not proof that the bytecode has no source or no behavior.



\subsection{RQ4: Source Reproduction}\label{sec:eval:runtime}

RQ4 reports whether selected source-recovery workflows can reproduce bytecode-level findings from RQ3.  We classify each finding by the first blocking point in the workflow: no selected decompiler available, decompilation failure, uncompilable source, behavior mismatch, or same-behavior reproduction.

\begin{table}[t]
  \centering
  \caption{Tool-bounded source reproduction for RQ3 findings. Replace placeholders with measured counts.}
  \label{tab:source-reproduction}
  \begin{tabular}{lrrrrr}
    \toprule
    Version & Findings & Decompiled & Compiled & Reproduced & Not reproduced \\
    \midrule
    3.10 & TBD & TBD & TBD & TBD & TBD \\
    3.11 & TBD & TBD & TBD & TBD & TBD \\
    3.12 & TBD & TBD & TBD & TBD & TBD \\
    \bottomrule
  \end{tabular}
\end{table}

For each abnormal outcome, we attempt decompilation, recompilation, and re-execution under the same runtime using the selected tools.  A finding is source-reproduced only if the recovered source produces the same behavior class.  Otherwise, it is classified as not reproduced by selected tools, with the concrete reason recorded as decompilation failed, recompilation failed, behavior changed, tool unavailable, or inconclusive.  This phrasing is deliberate: it avoids claiming that no source reproducer exists when the study only shows that the selected tools did not produce one.
